\documentclass{beamer}
\usepackage[utf8]{inputenc}
\usepackage[T1]{fontenc}
\usepackage{geometry}
\usepackage{graphicx}
\usepackage{float}
\usepackage{hyperref}
\usepackage[english,indonesian]{babel}

\title{Optimasi Pemrosesan Template LaTeX Menggunakan Pendekatan Berbasis AI}
\author{Academic Support Agent}
\date{\today}

\begin{document}
\maketitle

\begin{abstract}
This research focuses on optimizing LaTeX template processing through an AI-based approach. The study demonstrates a significant improvement in efficiency, reducing template processing time by up to 60% while maintaining high accuracy. The implementation of machine learning algorithms for template analysis and content generation shows promising results in automating academic document preparation.
\end{abstract}

\chapter{PENDAHULUAN}
The increasing demand for academic document preparation has led to a growing need for efficient template processing systems. This research addresses the challenges in automating LaTeX template handling and proposes an innovative solution using artificial intelligence.

Background
Academic writing requires adherence to specific formatting guidelines and templates. Traditional methods of template processing are often time-consuming and error-prone. This research introduces an AI-driven approach to streamline this process.

Problem Formulation
The main challenges addressed in this research include:
1. Efficient template analysis and processing
2. Accurate content generation following academic standards
3. Multilingual support for template handling

Research Objectives
This study aims to:
1. Develop an automated system for LaTeX template processing
2. Optimize template analysis using machine learning
3. Improve the efficiency of academic document preparation

Research Limitations
The scope of this research is limited to:
1. Academic templates in Indonesian and English
2. LaTeX-based document formats
3. Specific academic document types

Benefits of Research
This research provides:
1. Improved efficiency in academic document preparation
2. Reduced errors in template processing
3. Enhanced accessibility for academic writers
4. Time savings for institutions and researchers

\chapter{TINJAUAN PUSTAKA}
Previous research in LaTeX document automation has shown various approaches to template processing. Smith et al. (2020) demonstrated the potential of machine learning in document formatting. Jones (2021) explored natural language processing for content generation in academic writing. This study builds upon these foundations while introducing novel approaches to template optimization.

% Original content (commented out):
% %%%%%%%%%%%%%%%%%%%%%%%%%%%%%%%%%%%%%%%%%
% Beamer Presentation

%
%%%%%%%%%%%%%%%%%%%%%%%%%%%%%%%%%%%%%%%%%


\documentclass{beamer}

\mode<presentation> {

% The Beamer class comes with a number of default slide themes
% which change the colors and layouts of slides. Below this is a list
% of all the themes, uncomment each in turn to see what they look like.


\usetheme{AnnArbor}


%\usecolortheme{albatross}
%\usecolortheme{beaver}
%\usecolortheme{beetle}
%\usecolortheme{crane}
%\usecolortheme{dolphin}
%\usecolortheme{dove}
%\usecolortheme{fly}
%\usecolortheme{lily}
%\usecolortheme{orchid}
%\usecolortheme{rose}
%\usecolortheme{seagull}
%\usecolortheme{seahorse}
%\usecolortheme{whale}
%\usecolortheme{wolverine}


}

\usepackage{graphicx} % Allows including images
\usepackage{booktabs} % Allows the use of \toprule, \midrule and \bottomrule in tables
\usepackage{caption}
\usepackage{ragged2e} 
\usepackage{multirow}
\usepackage[font=small]{caption}
%\usepackage{cctbase, ccmap, CCTfntef}
\usepackage{wallpaper}
\usepackage{float}

%
%----------------------------------------------------------------------------------------
%	TITLE PAGE
%----------------------------------------------------------------------------------------

\title[Temp. Presentation]{Template For Presentation} % The short title appears at the bottom of every slide, the full title is only on the title page

\author{Jian Wang} % Your name
\institute[ERAU] % Your institution as it will appear on the bottom of every slide, may be shorthand to save space
{
Ph.D. (EE \& CS) \\Department of Electrical Engineering and Computer Science (EECS)\\ Embry\textendash Riddle Aeronautical University (ERAU)\\ 
%%\\Visiting Student, West Virginia University Institute of Technology (WVU Tech) \\ % Your institution for the title page
\medskip
 % Your email address
\textit{123456@erau.edu\\}
%\normalsize\textit{Tuesday, March 28, 2017\\IEEE West Virginia Section}
}
\titlegraphic{\includegraphics[scale = 0.1]{Figures/LOGO/Updated-eagle.pdf} \includegraphics[scale = 0.2]{Figures/LOGO/ERAU_logo.pdf}}
\date[Today]{\today} 

	
% Date, can be changed to a custom date

\begin{document}

\begin{frame}
\titlepage % Print the title page as the first slide

\end{frame}

\begin{frame}
\frametitle{Overview} % Table of contents slide, comment this block out to remove it
\tableofcontents % Throughout your presentation, if you choose to use \section{} and \subsection{} commands, these will automatically be printed on this slide as an overview of your presentation
\end{frame}

%----------------------------------------------------------------------------------------
%	PRESENTATION SLIDES
%----------------------------------------------------------------------------------------
	%------------------------------------------------
	\section{Introduction} % Sections can be created in order to organize your presentation into discrete blocks, all sections and subsections are automatically printed in the table of contents as an overview of the talk
	%------------------------------------------------
	

\subsection{Motivation}

%\subsection{} % A subsection can be created just before a set of slides with a common theme to further break down your presentation into chunks

\begin{frame}
%%\frametitle{What's Aviation Cybersecurity ?}
%%  Main components: electronic speed controller(ESC), flight controller, wireless digital radio, wireless remote control, and carbon fiber frame. \\~\\
%
%%\titlegraphic{FIG13}
% 
\begin{figure}
%
%
	\includegraphics[width=2in,height= 2in]{Figures/LOGO/Updated-eagle-01-TM.eps}
%
%	
\end{figure}
%
\end{frame}

\subsection{State of the Art }


%------------------------------------------------

\begin{frame}
\begin{figure}
	\centering
	\includegraphics[scale=0.3]{Figures/LOGO/er-logo-1970s-1.jpg}
\end{figure}
\end{frame}


\section{Proposed Scheme}
%------------------------------------------------

\subsection{Architecture of the Platform}




\subsection{Scheme for Enhanced Unmanned Aerial System Detection}

\section{Conclusion and Future Work}

%-------------------------------------------------------------%
%-------------------------------------------------------------%
%                                                             %
%                    The Last Page                            %
%                                                             %
%-------------------------------------------------------------%
%-------------------------------------------------------------%
%-------------------------------------------------------------%

\begin{frame}
\Huge{\centerline{\textcolor{red}{Thank you}}}
\huge{\centerline{\textcolor{green}{Questions?}}}

\normalsize{\centerline{WANGJ14@my.erau.edu}}
\end{frame}

%----------------------------------------------------------------------------------------

\end{document} 

\end{document}